\subsection{Dataset}

We use the \emph{Electric Vehicle Charging Station Data} dataset
published by the City of Boulder, Colorado on its Open Data
portal~\cite{boulderdata}. The dataset is updated regularly; the
snapshot we work with covers the period from August 22, 2018 to
November 30, 2023 (1\,927 days), and contains 148\,136 individual
charging sessions across 26 city-owned level-2 stations. Each session
record carries 15 fields, of which the most relevant for forecasting
are the station identifier, the start and end timestamps, the
charging duration, and the energy delivered in kWh. There are
essentially no missing values: only four sessions out of 148\,136 have
a null end-time.

We focus on the station with the highest session volume,
\textbf{N~Boulder Rec~1}, with 16\,278 sessions in the period of
interest. The reason for picking a single station is simple: with one
station the daily series is dense and continuous, with only 5.2\% of
days at zero. Aggregating across all 26 stations would smooth out the
station-level structure we want to learn (commute patterns at
recreational vs.\ commercial locations are very different), and
working with multiple separate series would require multi-task models
that go beyond the scope of this project. The sessions for
N~Boulder Rec~1 are summed by date to produce a univariate daily
series of 1\,927 observations, with mean 74.6~kWh/day, standard
deviation 50.3~kWh, and a range of 0--310~kWh. 100 days have zero
energy (mostly holidays or short station outages); we preserve these
zeros instead of imputing them, since they reflect real operational
patterns that the model should learn.

\begin{figure}[!htbp]
\centering
\includegraphics[width=0.95\textwidth]{fig_series}
\caption{Daily energy consumption at N~Boulder Rec~1 from August 2018
to November 2023. The orange band marks the 30-day held-out test
window; everything to the left of it is training data. The COVID-19
drop in 2020 and the post-pandemic recovery are clearly visible.}
\label{fig:series}
\end{figure}

\subsection{Exogenous Variables}

For the models that support them (SARIMAX, ARIMAX, LSTM, PatchTST), we
add three exogenous channels aligned with the daily target:
\begin{itemize}[leftmargin=*]
    \item \textbf{Temperature.} Daily mean temperature in Boulder
    pulled from the Open-Meteo historical weather API. We picked mean
    over min/max because most charging happens during commute hours,
    when the mean is the closest proxy for the temperature actually
    seen by the chargers.
    \item \textbf{Day-of-week flags.} A binary \emph{is\_weekend}
    flag, plus the integer day of week as a feature when relevant. We
    do \emph{not} use a holiday-aware day-of-week encoding because
    holidays are handled separately.
    \item \textbf{Holiday indicator.} A binary flag for US public
    holidays, generated from the \texttt{holidays} Python package
    using the federal calendar plus Colorado state holidays.
\end{itemize}
These three variables are known at forecast time (day-of-week and
holiday from the calendar, temperature from a standard one-week
weather forecast), so feeding them as future exogenous inputs is
causally legitimate.

\subsection{Stationarity and Seasonality}

We ran the Augmented Dickey-Fuller test on the raw daily series. The
ADF statistic is $-3.097$ with $p$-value $0.027$ at 24 lags (AIC
selection), which lets us reject the unit-root null at the 5\% level.
The series is stationary, not a random walk; we therefore use $d=0$
as the default differencing order for the (S)ARIMA family, with the
search space allowing $d \in \{0,1\}$ for SARIMA's first difference of
the seasonally differenced series.

Figure~\ref{fig:acf} shows the autocorrelation function of the raw
series. The picture is unambiguous: clear peaks at lags 7, 14, 21, 28
and 35, with a slow decay in between, telling us the weekly cycle is
the dominant seasonal component. There is no annual peak visible in
this lag range because we only show 42 lags, but the seasonal
behaviour at the year scale is clear in the time-series plot
(Fig.~\ref{fig:series}, where summer months consistently sit above
the mean). The PACF (not shown for space reasons) has significant
spikes at lags 1 and 7, suggesting an AR(1) main component with a
seasonal AR at lag 7 --- which is the search shape we use for the
(S)ARIMA grid.

\begin{figure}[!htbp]
\centering
\includegraphics[width=0.85\textwidth]{fig_acf}
\caption{Autocorrelation of the daily series. The repeating peaks at
multiples of 7 are the signature of a strong weekly seasonal cycle;
this is the structural reason classical (S)ARIMA models do not need
to be replaced by anything fancier to be competitive.}
\label{fig:acf}
\end{figure}

\subsection{Train / Test Split and Evaluation Protocol}

The last 30 days of the series (Nov 1, 2023 -- Nov 30, 2023) are held
out as the test set; everything earlier is training data. Within the
training data, the last 60 days are reserved as a validation split,
used by the LSTM and PatchTST scripts for early stopping. Total
\textbf{training data: 1\,897 days}; \textbf{validation: 60 days
inside training}; \textbf{test: 30 days}.

We use rolling-origin (also known as walk-forward) evaluation with
\textbf{four non-overlapping windows of seven days each}. At window
$i$ the model sees the full training set plus the first
$i \cdot 7$ days of the test set as history, and produces a 7-day
forecast for days $i \cdot 7 + 1, \dots, (i+1) \cdot 7$ of the test
set. Metrics are computed on each window separately and then averaged.
This protocol mimics what a real operator would do: every Monday they
have access to the previous week's actuals and need to forecast the
upcoming week. With 30 test days we get four such weekly windows; the
remaining two days are dropped (we keep windows non-overlapping so
metrics across windows are independent).

\subsection{Metrics}

We report seven metrics per window and aggregate them as
mean~$\pm$~standard deviation across the four windows:
\begin{itemize}[leftmargin=*]
    \item \textbf{MAE} (mean absolute error) in kWh --- our primary
    metric, because it is in the same units as the target and easy to
    interpret operationally (``the model is off by 25 kWh on
    average''). The standard deviation across windows is at least as
    important as the mean: a model with low mean but high variance is
    risky in production.
    \item \textbf{RMSE} (root mean squared error) in kWh, which
    penalises large errors more than MAE.
    \item \textbf{MAPE} and \textbf{sMAPE} (mean absolute percentage
    error and its symmetric variant). MAPE blows up on zero actuals
    so we report both.
    \item \textbf{MASE} (mean absolute scaled error), scaled by the
    in-sample MAE of the seasonal-naive forecast. MASE $< 1$ means
    the model beats the trivial baseline; MASE $= 0.65$ means it cuts
    that error by 35\%.
    \item \textbf{Bias} (mean signed error) in kWh --- positive means
    the model over-predicts on average. Useful for spotting
    systematic offsets.
    \item \textbf{R\textsuperscript{2}} on the test window. R\textsuperscript{2}
    can be negative when the model is worse than the mean of the
    test window, which is common with very short windows; we still
    report it for completeness.
\end{itemize}
Throughout the paper we rank by MAE because it is the most operationally
meaningful metric. The full table with every metric is in the project
repository.

\subsection{Reproducibility and Heavy-Training Scripts}

All the heavy training lives in standalone Python scripts:
\texttt{scripts/sarima\_grid\_search.py} for the AIC grid,
\texttt{scripts/lstm\_train.py} for the LSTM sweep, and
\texttt{scripts/patchtst\_train.py} for the PatchTST seed ensemble.
The notebooks load the resulting cached predictions and apply the
shared \texttt{scripts/evaluation.py} rolling-origin routine, so a
single notebook never has to retrain a transformer to produce its
leaderboard row. Training the PatchTST seed ensemble takes
$\approx$2~min per seed on an RTX 5070 Ti.
